% 第6章 6.6：谱响应曲线。
% (a) 三点路径谱 {0,1,3}：目标 g(l)=1/(1+l)；在三个特征值上插值的二次多项式
%     p(l)=1-5l/8+l^2/8；一次多项式 h(l)=1-l/3（响应 1,2/3,0）。
% (b) 区间 [0,2] 上目标 g(l)=1/(1+3l) 的 Chebyshev 节点插值多项式，
%     系数由同目录 chebyshev_coefficients.py 计算：K=2 最大误差 0.1824，K=4 最大误差 0.0375。
\begin{tikzpicture}[font=\figurefont\footnotesize]
\pgfplotsset{every axis/.style={width=54mm,height=38mm,axis lines=left,
  x axis line style={draw=FigureStroke,line width=.6pt},y axis line style={draw=FigureStroke,line width=.6pt},
  tick style={draw=FigureStroke},tick label style={font=\figurefont\scriptsize,text=FigureInk},
  label style={font=\figurefont\scriptsize},xlabel={$\lambda$},
  xlabel style={at={(axis description cs:1,0)},anchor=north west,yshift=1mm},
  ymajorgrids,grid style={draw=FigureGrid,line width=.4pt},
  title style={font=\figurefont\footnotesize},
  legend style={draw=none,fill=none,font=\figurefont\scriptsize,cells={anchor=west}},
  samples=120}}
\begin{axis}[name=A,xmin=0,xmax=3.3,ymin=-0.15,ymax=1.1,xtick={0,1,2,3},ytick={0,0.5,1},title={(a) 三点路径，谱为$\{0,1,3\}$},legend style={at={(0.5,-0.24)},anchor=north}]
\addplot[draw=FigureBlue,line width=1.1pt,domain=0:3.3] {1/(1+x)};
\addplot[draw=FigureGreen,line width=1.1pt,dashed,domain=0:3.3] {1-5*x/8+x^2/8};
\addplot[draw=FigureYellow,line width=1.1pt,dash dot,domain=0:3.3] {1-x/3};
\addplot[only marks,mark=*,mark size=1.4pt,FigureBlue] coordinates {(0,1) (1,0.5) (3,0.25)};
\addplot[only marks,mark=square*,mark size=1.2pt,FigureYellow] coordinates {(0,1) (1,0.6667) (3,0)};
\legend{$g(\lambda)=1/(1+\lambda)$,$p(\lambda)$，插值$g$,$h(\lambda)=1-\lambda/3$}
\end{axis}
\begin{axis}[at={(A.east)},anchor=west,xshift=10mm,xmin=0,xmax=2.05,ymin=-0.15,ymax=1.1,xtick={0,0.5,1,1.5,2},ytick={0,0.5,1},title={(b) 区间$[0,2]$上的多项式逼近},legend style={at={(0.5,-0.24)},anchor=north}]
\addplot[draw=FigureBlue,line width=1.1pt,domain=0:2] {1/(1+3*x)};
\addplot[draw=FigureYellow,line width=1.1pt,dash dot,domain=0:2] {0.817568-0.810811*x+0.243243*x^2};
\addplot[draw=FigureGreen,line width=1.1pt,dashed,domain=0:2] {0.962523-1.950648*x+2.104257*x^2-1.066009*x^3+0.199877*x^4};
\legend{$1/(1+3\lambda)$,$K=2$,$K=4$}
\end{axis}
\end{tikzpicture}
